Devido ao seu alto desempenho e à sua natureza como linguagem compilada de baixo 
nível, a linguagem C apresenta tanto vantagens quanto limitações. Entre suas 
principais vantagens, destacam-se:

\begin{itemize}
        \renewcommand{\labelitemi}{\tiny$\blacksquare$}
    \item Alta portabilidade entre diferentes arquiteturas e sistemas
    \item Alto desempenho comparado a outras linguagens
    \item Controle direto e eficiente sobre o uso de memória
    \item Acesso direto a recursos do sistema e ao hardware
\end{itemize}

Entretanto, essas vantagens vêm acompanhadas de um custo. Essa ideia pode ser 
substanciada por uma frase que gosto bastante:

\begin{quote}
"No começo todos queremos resultados, mas no fim, tudo que queremos é controle"
\footnote{Eskil Steenberg}
\end{quote}

No cenário atual da indústria de software, frequentemente orientado por prazos curtos 
e rápida entrega de funcionalidades, observa-se a crescente adoção de linguagens de 
alto nível. Essas linguagens permitem o desenvolvimento mais ágil, muitas vezes à 
custa de desempenho, previsibilidade e controle fino sobre os recursos do sistema.

Nesse contexto, a linguagem C se destaca por sua simplicidade estrutural e pelo baixo 
nível de abstração. O que acaba transferindo ao programador a responsabilidade 
por diversos aspectos que, em linguagens mais modernas, são automatizados ou abstraídos 
por bibliotecas e frameworks, como gerenciamento de memória, manipulação de estruturas 
de dados e controle explícito de recursos.

Apesar dessa exigência adicional, o domínio da linguagem C proporciona uma compreensão 
profunda dos fundamentos da computação, incluindo organização de memória, fluxo de 
execução e interação com o hardware. Esse conhecimento serve como base sólida para o 
aprendizado de outras linguagens e paradigmas.

Além disso, a linguagem C continua sendo amplamente utilizada em áreas onde desempenho, 
previsibilidade e controle são essenciais, como:

\begin{itemize}
        \renewcommand{\labelitemi}{\tiny$\blacksquare$}
    \item Desenvolvimento de sistemas operacionais
    \item Programação com sistemas embarcados e microcontroladores
    \item Implementação de drivers de dispositivos
    \item Construção de compiladores e outras linguagens de programação
    \item Criação de utilitários e ferramentas do sistema
\end{itemize}
